\documentclass[UTF8,11pt,fontset=none]{ctexart}
\usepackage[a4paper,margin=25mm,headheight=15pt,footskip=13mm]{geometry}
\usepackage{fontspec}
\IfFontExistsTF{TeX Gyre Termes}{\setmainfont{TeX Gyre Termes}}{\setmainfont{TeXGyreTermes}}
\IfFontExistsTF{TeX Gyre Heros}{\setsansfont{TeX Gyre Heros}}{\setsansfont{TeXGyreHeros}}
\IfFontExistsTF{DejaVu Sans Mono}{\setmonofont{DejaVu Sans Mono}[Scale=MatchLowercase]}{\setmonofont{Latin Modern Mono}}
\IfFontExistsTF{Noto Serif CJK SC}{\setCJKmainfont{Noto Serif CJK SC}[AutoFakeSlant=0.15]}{\setCJKmainfont{FandolSong-Regular.otf}[BoldFont=FandolSong-Bold.otf,ItalicFont=FandolKai-Regular.otf]}
\IfFontExistsTF{Noto Sans CJK SC}{\setCJKsansfont{Noto Sans CJK SC}}{\setCJKsansfont{FandolHei-Regular.otf}}
\IfFontExistsTF{Noto Sans Mono CJK SC}{\setCJKmonofont{Noto Sans Mono CJK SC}}{\setCJKmonofont{FandolFang-Regular.otf}}
\usepackage{amsmath,amssymb,amsthm,mathtools,mathrsfs,esint,bm}
\usepackage{booktabs,longtable,array,calc}
\usepackage{enumitem}
\setlist{nosep,leftmargin=1.7em,topsep=4pt}
\usepackage{fvextra}
\DefineVerbatimEnvironment{Highlighting}{Verbatim}{breaklines,breakanywhere,fontsize=\small}
\usepackage{xurl}
\usepackage[unicode,hidelinks,bookmarksnumbered]{hyperref}
\usepackage{bookmark}
\usepackage{fancyhdr}
\pagestyle{fancy}
\fancyhf{}
\fancyhead[L]{\small 第一科学主线：近期数学补充}
\fancyhead[R]{\small 2026-10-07}
\fancyfoot[C]{\thepage}
\renewcommand{\headrulewidth}{0.3pt}
\providecommand{\tightlist}{\setlength{\itemsep}{0pt}\setlength{\parskip}{0pt}}
\providecommand{\passthrough}[1]{#1}
\setlength{\emergencystretch}{3em}
\setlength{\parindent}{2em}
\setlength{\parskip}{3.5pt plus 1pt}
\linespread{1.12}
\setlength{\abovedisplayskip}{10pt plus 2pt minus 2pt}
\setlength{\belowdisplayskip}{10pt plus 2pt minus 2pt}
\allowdisplaybreaks[2]
\renewcommand{\arraystretch}{1.15}
\let\cleardoublepage\clearpage
\setcounter{tocdepth}{1}
\setcounter{secnumdepth}{2}
\newtheorem{theorem}{定理}[section]
\newtheorem{proposition}[theorem]{命题}
\newtheorem{lemma}[theorem]{引理}
\newtheorem{corollary}[theorem]{推论}
\theoremstyle{definition}
\newtheorem{definition}[theorem]{定义}
\newtheorem{assumption}[theorem]{前置}
\theoremstyle{remark}
\newtheorem{remark}[theorem]{说明}
\newcommand{\Ric}{\operatorname{Ric}}
\newcommand{\Vol}{\operatorname{Vol}}
\newcommand{\Var}{\operatorname{Var}}
\newcommand{\Rm}{\operatorname{Rm}}
\newcommand{\supp}{\operatorname{supp}}
\newcommand{\tr}{\operatorname{tr}}
\newcommand{\id}{\operatorname{Id}}
\newcommand{\dd}{\,\mathrm d}
\pdfstringdefDisableCommands{\def\fint{int}\def\mathcal#1{#1}\def\mathfrak#1{#1}\def\mathscr#1{#1}\def\operatorname#1{#1}\def\texorpdfstring#1#2{#2}}

\renewcommand{\thesection}{3.\arabic{section}}
\title{第一科学主线：近期研究补充\\[6pt]\large 3．全度量 Ricci 间隙的定量动态后果}
\author{数学研究交接稿}\date{2026年10月7日}
\begin{document}
\maketitle
\tableofcontents
\clearpage
本分卷保留续研中给出的完整组合证明。前提是一个真正的全度量正 \(P\)-gap、负 Yamabe 不变量，以及指定外部四维 Ricci 流延拓定理。它是下游应用，不被登记为无密度 O10 或 O11 最优首项的解答。来源中的推导在下文完整保留。\cite{R02,OAI351Flow,Ishida2008}

\section{统一的静态输入}\label{v3-ux7edfux4e00ux7684ux9759ux6001ux8f93ux5165}

令 M 为闭连通光滑四流形，设

\[
 P(g)=\int_M|\operatorname{Ric}^{\circ}_g|^2\,dV_g\ge p_0>0
 \quad\text{对 M 上每一个光滑度量成立},
 \qquad \mathcal Y(M)<0.
\]

这两项输入必须在对应实际供给范围内验证。只有 \(P+D_2\) 的正下界不够。

第一主线的一个直接输入方式是：存在统一 \(A>0\) 与 \(j_0>0\)，使每个竞争度量满足

\[
 \mathcal J(g)=\frac{4P(g)}{R(g)}\ge j_0,
 \qquad S_2(g)=\int s_g^2\ge32\pi^2A,
 \qquad R(g)=P(g)+\frac14S_2(g).
\]

于是可以取

\[
 \boxed{p_0=2\pi^2 A j_0.}
\]

例如从 N1 取得 \(j_0\) 时，必须保留拓扑超额项：\(\mathcal J+d\ge\varepsilon\) 只有在 \(\varepsilon-d>0\) 时给出正 \(j_0\)。短缺版同理。小 h 平衡族则使用已经证明的线性正 \(\mathcal J\)-下界。

若同一全度量正投影供给还给出负 Yamabe 上界

\[
 \mathcal Y(M)\le-4\pi\sqrt{2A},
\]

则下文可以取 \(a=4\pi\sqrt{2A/V}\)。这可由各个负 Yamabe 代表上的标量预算得到，并非只从任意一个度量的标量积分得到。

\section{外部输入：四维有限时间标量延拓}\label{v3-ux5916ux90e8ux8f93ux5165ux56dbux7ef4ux6709ux9650ux65f6ux95f4ux6807ux91cfux5ef6ux62d3}

采用 openai/math 第351族中的下述结果。

\textbf{外部定理 F。} 若闭光滑四流形上的未归一化 Ricci 流 \(\partial_\tau\widetilde g=-2\operatorname{Ric}(\widetilde g)\) 在有限区间 \([0,\widetilde T)\) 满足 \(\sup|s_{\widetilde g}|<\infty\)，则它在原光滑流形上可以光滑延拓。

精确源文件： \path{preprints/Bounded-scalar-curvature-and-smooth-extension-of-four-dimensional-Ricci-flow-September-24-2026/build/sections/01-introduction.tex}，\texttt{thm:extension}。

这里借用外部定理的实际结论，不把其证明中要求的 Ricci 流时空估计改名成静态 O10 的已知预算。

\section{定量阈值定理}\label{v3-ux5b9aux91cfux9608ux503cux5b9aux7406}

考虑体积归一化 Ricci 流

\[
 \partial_t g=-2\operatorname{Ric}+\frac12\bar s\,g,
 \qquad \bar s(t)=\fint_M s_{g(t)},
 \qquad V=\operatorname{Vol}_{g(0)}M.
\]

定义初始数据

\[
 m_0=\min_Ms_{g(0)},\qquad b_0=-m_0>0,
 \qquad \bar s_0=\fint_Ms_{g(0)}.
\]

任选 \(a>0\) 使

\[
 \mathcal Y(M)\le-a\sqrt V.
\]

则 \(b_0\ge a\)。对任意 \(K>\max\{0,\max_Ms_{g(0)}\}\)，令

\[
 \tau_K=\inf\{t:\max_Ms_{g(t)}\ge K\}.
\]

\textbf{定理 A（原续研组合推导）。} 在第1节静态输入及外部定理 F 下，\(\tau_K\) 在最大光滑时间之前存在，且

\[
 \boxed{
 \tau_K\le\frac{V}{2p_0}
 \left[
 K\left(1+\log\frac{b_0}{a}\right)-\bar s_0+b_0-a
 \right].
 }
\tag{3.A}
\]

它既不预设流永生，也不预设曲率、直径或 Sobolev 常数统一有界。

\subsection{Yamabe 上界约束标量最小值}\label{v3-yamabe-ux4e0aux754cux7ea6ux675fux6807ux91cfux6700ux5c0fux503c}

令 \(m(t)=\min s_{g(t)}\)。由于 Yamabe 不变量为负，每个时刻有 \(m(t)<0\)。对于任意非零正函数 u，

\[
 \frac{\int(6|du|^2+s u^2)}{(\int u^4)^{1/2}}
 \ge m(t)\frac{\int u^2}{(\int u^4)^{1/2}}
 \ge m(t)\sqrt V.
\]

最后一个方向使用 \(m(t)<0\) 与 Hölder。取下确界并使用 \(Y([g(t)])\le\mathcal Y(M)\)，得

\[
 \boxed{m(t)\le-a<0.}
\tag{3.3.1}
\]

\subsection{两条精确演化关系}\label{v3-ux4e24ux6761ux7cbeux786eux6f14ux5316ux5173ux7cfb}

使用 \(\Delta=\operatorname{div}\nabla\)，标量演化为

\[
 \partial_t s=\Delta s+2|\operatorname{Ric}^{\circ}|^2
             +\frac12s(s-\bar s).
\]

最小值是局部 Lipschitz 函数；在其几乎处处可微处，最大值原理给

\[
 m'\ge\frac{-m}{2}(\bar s-m)\ge0.
\tag{3.3.2}
\]

因此 \(m_0\le m(t)\le-a\)。

另一方面，\(\partial_t dV=(\bar s-s)dV\)，故

\[
 \boxed{
 \bar s'=\frac{2P(g(t))}{V}-\frac12\operatorname{Var}(s),
 \quad \operatorname{Var}(s)=\fint(s-\bar s)^2.
 }
\tag{3.3.3}
\]

\subsection{有界标量区间最多能持续多久}\label{v3-ux6709ux754cux6807ux91cfux533aux95f4ux6700ux591aux80fdux6301ux7eedux591aux4e45}

在任意满足 \(s\le K\)、\(K\ge0\) 的光滑时间区间 \([0,t]\) 内，点态 \((s-m)(K-s)\ge0\) 给

\[
 \operatorname{Var}(s)\le(K-\bar s)(\bar s-m)
                      \le(K-m)(\bar s-m).
\]

由式(3.3.2)，

\[
 \frac12\operatorname{Var}(s)
 \le\left(1+\frac{K}{-m}\right)m'.
\]

令 \(u_t=-m(t)\)，积分得

\[
 \frac12\int_0^t\operatorname{Var}(s)\,du
 \le K\log\frac{b_0}{u_t}+b_0-u_t
 \le K\log\frac{b_0}{a}+b_0-a.
\tag{3.3.4}
\]

这里最后一步使用 \(u_t\ge a\) 和 \(K\ge0\)。将式(3.3.4)与(3.3.3)积分合并，并使用 \(P\ge p_0\)、\(\bar s(t)\le K\)，得到

\[
 \boxed{
 \frac{2p_0}{V}t
 \le K\left(1+\log\frac{b_0}{a}\right)-\bar s_0+b_0-a.
 }
\tag{3.3.5}
\]

这条不等式本身不使用新延拓定理。

\subsection{新延拓定理排除提前终止的分支}\label{v3-ux65b0ux5ef6ux62d3ux5b9aux7406ux6392ux9664ux63d0ux524dux7ec8ux6b62ux7684ux5206ux652f}

若流在达到 K 之前已有有限最大时间 T，则在 \([0,T)\) 上 \(m_0\le s\le K\)。因此 \(\bar s\) 有界。

定义

\[
 c(t)=\exp\left(-\frac12\int_0^t\bar s(u)du\right),
 \quad \widetilde t=\int_0^t c(u)du,
 \quad \widetilde g(\widetilde t)=c(t)g(t).
\]

在有限区间上 c 有正的上下界。直接求导可核对 \(\partial_{\widetilde t}\widetilde g=-2\operatorname{Ric}(\widetilde g)\)， 且 \(s_{\widetilde g}=s_g/c\) 有界。外部定理 F 允许延拓 \(\widetilde g\)，再按其体积归一化即可延拓 g，矛盾。

若永远没有达到 K，且流永生，则式(3.3.5)与 \(t\to\infty\) 矛盾。因此 \(\tau_K<T_{\max}\) 存在；在 \(t<\tau_K\) 使用(3.3.5)，令 \(t\uparrow\tau_K\)，得到(3.A)。证毕。

\section{两项直接后果}\label{v3-ux4e24ux9879ux76f4ux63a5ux540eux679c}

\subsection{永生时，历史最大正标量至少线性增长}\label{v3-ux6c38ux751fux65f6ux5386ux53f2ux6700ux5927ux6b63ux6807ux91cfux81f3ux5c11ux7ebfux6027ux589eux957f}

定义

\[
 K(t)=\max\left\{0,\sup_{0\le u\le t}\max_Ms_{g(u)}\right\}.
\]

将 K(t) 用作式(3.3.5)的区间上界，有

\[
 \boxed{
 K(t)\ge
 \frac{(2p_0/V)t+\bar s_0-b_0+a}
 {1+\log(b_0/a)}.
 }
\tag{3.4.1}
\]

因而若流永生，

\[
 \liminf_{t\to\infty}\frac{K(t)}t
 \ge\frac{2p_0/V}{1+\log(b_0/a)}>0.
\]

这是历史最大值的增长，不是每个时刻 \(\max_Ms_{g(t)}\) 都单调，也不是有限时间爆破率。

\subsection{第一主线的显式拓扑版本}\label{v3-ux7b2cux4e00ux4e3bux7ebfux7684ux663eux5f0fux62d3ux6251ux7248ux672c}

若使用 \(p_0=2\pi^2Aj_0\)、\(a=4\pi\sqrt{2A/V}\)，则

\[
 \boxed{
 \tau_K\le\frac{V}{4\pi^2Aj_0}
 \left[K\left(1+\log\frac{b_0}{a}\right)-\bar s_0+b_0-a\right].
 }
\tag{3.4.2}
\]

若初始度量已归一化为 \(s=-4\)，则 \(b_0=4\)、\(\bar s_0=-4\)、\(a\le4\)，从而

\[
 \boxed{
 \tau_K\le\frac{V}{2p_0}
 \left[K\left(1+\log\frac4a\right)+8-a\right].
 }
\tag{3.4.3}
\]

对原平衡族，已经成立的线性 \(\mathcal J\)-gap 给 \(p_{0,j}\ge cA_jh_j\)。对选定的 \(s=-4\) 初始代表，若已有 \(V_j\le C A_j\)，则 a 有统一正下界，故

\[
 \boxed{\tau_{K,j}\le\frac{C_K}{h_j}.}
\tag{3.4.4}
\]

原稿可选的低缺陷代表满足这一体积预算：由 \(8\pi^2A+P\le R\le8\pi^2AU+4P\)， \(U=1+O(h)\)、\(\mathcal J=O(h)\)，得 \(R=O(3.A)\)，再用 \(s=-4\) 的 \(V=(R-P)/4\)。

式(3.4.4)是上界，不是精确寿命公式，也不决定 O11 的系数。它使用已证线性阶，不需要先完成无密度 O10 或 O11。

\begin{thebibliography}{E99}
\small
\bibitem[R02]{R02} 本会话证明补充，\emph{新外部定理的实际接入与定量 Ricci 流推论}，2026-10-07。\path{source_notes/R02_Quantitative_Ricci_Flow.md}。

\bibitem[E03]{OAI351Flow} OpenAI，\emph{Bounded scalar curvature and smooth extension of four-dimensional Ricci flow}，2026-09-24。同一指定版本目录：\path{preprints/Bounded-scalar-curvature-and-smooth-extension-of-four-dimensional-Ricci-flow-September-24-2026/}；主定理 \path{thm:extension}。

\bibitem[E05]{Ishida2008} M. Ishida，\emph{The Normalized Ricci Flow on Four-Manifolds and Exotic Smooth Structures}，arXiv:0807.2169，§§2--3。\url{https://arxiv.org/abs/0807.2169}。

\end{thebibliography}

\end{document}
